\section{Theory and numerical support}\label{app:theory-support}

The equilibrium proof combines the curvature of proportional shares with
strictly curved capacity costs; the response proof then restricts that field
to the locally unchanged active coordinates.
The composition and cash-flow arguments identify which parts of the resulting
capacity response survive normalization and comparisons across firms.
Numerical certificates bound operating-solver error before article sampling.

\subsection{Contest inheritance and the capacity
adaptation}\label{app:contest-inheritance}

The economic technology in \eqref{eq:firm-problem} adapts the multi-contest
framework of \citet{xu_zenou_zhou_conflict_2022}.
Their Equations~1--2 allocate effort directly to contests, while Remark~5,
Equation~54 and its footnote allow a positive, contest-specific denominator
term.
Setting contest production to the identity and $K=I$ recovers the corresponding
regularized share rule.
Here $Kq_i$ instead maps operating capacities to appeal across demands, and
outside appeal remains positive rather than being taken to zero in an
existence approximation.

Their Proposition~1 characterizes Nash equilibrium by a variational
inequality; Propositions~2--3 and Theorem~2 use monotonicity for uniqueness.
Theorem~3 derives local active-coordinate responses in their pure-cost setting
under their cost assumptions and nondegeneracy,
which corresponds to the strict complementarity condition used here.

The capacity adaptation requires its own hypotheses.
Xu and colleagues maintain increasing convex costs and impose strong increase
where needed; Assumption~\ref{ass:primitives} instead uses uniformly positive
quadratic curvature, including a general symmetric cost matrix.
Positive definiteness alone does not imply coordinatewise increasing cost.
Although one may represent feasible appeals as
$\mathcal A_i=\{Kq_i:q_i\geq0\}$ and minimize capacity cost over each appeal's
preimage, the resulting effective cost need not inherit their smoothness and
monotonicity assumptions.
Their restricted-contest Proposition~5 therefore does not make the present
orthant result an automatic corollary.
The following proofs establish the required concavity, monotonicity,
existence and sensitivity for this specification.

\subsection{Anchored operating stages}\label{app:anchored-stage}

For a fixed preceding allocation $a$, the stage field is
$F(q;r)+\kappa(q-a)$.
Its difference between any two feasible profiles has inner product at least
$(\nu+\kappa)\|q-q'\|^2$ with $q-q'$, because the anchor cancels and the
static field is $\nu$-strongly monotone.
The same coercivity argument as in Appendix~\ref{app:equilibrium} gives
existence on the capacity orthant; strong monotonicity gives uniqueness.
The shifted linear cost $c_i-\kappa a_i$ need not be nonnegative, so this
argument uses the anchored-stage result rather than imposing the static
nonnegative-cost restriction on the shifted problem.
Every stage retains contemporaneous rival choices and a fixed inherited
anchor, without a continuation-value first-order condition.

\subsection{Discrete-choice interpretation of proportional
appeal}\label{app:logit}

The interpretation in Remark~\ref{rem:logit} follows the conditional-logit
random-utility construction of
\citet{mcfadden_conditional_logit_1974} and
\citet{train_discrete_choice_2009}.
Fix a demand node $\ell$ and nonnegative integer counts $q_{ia}$.
Consumer choice is among all individual varieties and one outside option.
For variety $v$ belonging to firm $i$ at activity $a$, write deterministic
utility $V_v=-\tau_A\|x_a-z_\ell\|^2/d$; write $V_0=\log b_\ell$ for
outside utility.
The shocks $\epsilon_v$ are independent across alternatives with
$\Pr(\epsilon_v\leq e)=\exp\{-\exp(-e)\}$.
Independence and this distribution give
\begin{equation}\label{eq:logit-choice}
  \Pr(v\text{ chosen}\mid z_\ell)
  =\frac{\exp V_v}{\sum_w\exp V_w}.
\end{equation}
Summing this probability over firm $i$'s varieties yields
\begin{equation}\label{eq:logit-firm-share}
  \Pr(i\mid z_\ell)
  =\frac{\sum_a q_{ia}\exp\{-\tau_A\|x_a-z_\ell\|^2/d\}}
  {b_\ell+\sum_j\sum_a q_{ja}\exp\{-\tau_A\|x_a-z_\ell\|^2/d\}}
  =\frac{A_{i\ell}}{A_{i\ell}+B_{i\ell}}.
\end{equation}
For real nonnegative $q$, the same ratio defines a weighted-choice
extension: an active firm-activity alternative has deterministic utility
$\log q_{ia}+\log K_{\ell a}$, and zero-weight alternatives are absent.
This extension retains the share formula without interpreting a fractional
mass as a fractional number of independent shocks.
Multiproduct aggregation is also central to
\citet{nocke_schutz_aggregative_2018}, whose endogenous-price model
has different strategic choices from this fixed-margin capacity game.

\begin{remark}[Logit interpretation]\label{rem:logit}
  The share rule is multinomial logit \citep{mcfadden_conditional_logit_1974}
  when $q_{ia}$ measures varieties, outlets, or participation and $c_{ia}$
  is linear cost per unit.
  The parameter $\tau_A$ measures mismatch cost relative to taste heterogeneity,
  and $r_\ell=M_\ell m_\ell$ combines market size with fixed per-sale margin.
  Conditional on other primitives, size and margin changes are
  observationally equivalent in firm cash flows.
\end{remark}

\subsection{Mathematical preliminaries and notation}\label{app:proofs}

All spaces in the economic model are finite dimensional.
Let $I=\{1,\ldots,N\}$, $\mathcal A=\{1,\ldots,m\}$, and
$\mathcal L=\{1,\ldots,L\}$ index firms, activities, and demand
nodes, with $N,m,L\geq1$.
Firm $i$ chooses $q_i\in\mathcal Q_i=\mathbb R_+^m$ and the joint
profile belongs to
$\mathcal Q=\prod_i\mathcal Q_i$.
We use the ordinary Euclidean product structure:
\begin{equation}\label{eq:app-inner}
  \langle q,v\rangle=\sum_i\sum_a q_{ia}v_{ia},\qquad
  \|q\|^2=\sum_i\|q_i\|^2.
\end{equation}
The cone $\mathcal Q$ is convex, so segments and one-sided deviations remain
feasible even at a boundary allocation.

The appeal matrix $K\in\mathbb R_+^{L\times m}$ is common to firms.
Its row $k_\ell^\top$ maps a capacity block into the scalar appeal
$A_{i\ell}=k_\ell^\top q_i$.
Put $B_{i\ell}=b_\ell+\sum_{j\ne i}A_{j\ell}$ and
$D_\ell=b_\ell+\sum_i A_{i\ell}=A_{i\ell}+B_{i\ell}$.
Outside appeal $b_\ell>0$ guarantees $D_\ell\geq b_\ell$ at every
feasible profile.
Opportunity weights $r_\ell\geq0$ are in cash-flow units, and
$R_\Sigma=\sum_\ell r_\ell$ is the total available opportunity.
No rank condition on $K$ is required: different capacity changes can
have the same appeal image.

Costs are $C_i(q_i)=c_i^\top q_i+q_i^\top H_i q_i/2$, with $c_i\geq0$,
$H_i=H_i^\top$, and $H_i\succeq\nu I$ for a common $\nu>0$.
The Gaussian benchmark economy uses the same $H$ across firms and
heterogeneous $c_i$.
The proofs permit different $H_i$ when the same lower bound holds.
The negative marginal-profit field is $F_i=-\nabla_i\pi_i$; its
components have units of
cash flow per unit capacity.
Cost curvature has units of cash flow per squared unit capacity.
Consequently the error and regret bounds in
Appendix~\ref{app:certificates} have capacity
and payoff units, respectively.

\subsection{Proportional-share payoff algebra}\label{app:payoff}

The margin component of the operating payoff and the individual capacity
choice are, in full notation,
\begin{equation}\label{eq:revenue}
  \mathcal M_i(q_i;r_t,B_i)=\sum_\ell r_{\ell t}
  \frac{A_{i\ell}(q_i)}{A_{i\ell}(q_i)+B_{i\ell}}.
\end{equation}
\begin{equation}\label{eq:firm-choice}
  q_i^*(t)=\arg\max_{q_i\geq0}\pi_i(q_i,q_{-i};r_t).
\end{equation}

For a unilateral deviation, fix rivals and one demand node, writing current
own appeal as $a\geq0$, rival-plus-outside appeal as $B>0$ and $D=a+B$.
A feasible appeal change $s$ can have either sign, but $a+s\geq0$ keeps both
denominators strictly positive.
Subtracting the share margins gives the exact identity
\begin{align}
  \frac{r(a+s)}{D+s}-\frac{ra}{D}
  &=\frac{r[(a+s)D-a(D+s)]}{D(D+s)}
  =\frac{rBs}{D(D+s)}\label{eq:app-share}\\
  &=\frac{rBs}{D^2}-\frac{rBs^2}{D^2(D+s)}.\label{eq:share-expansion}
\end{align}
The nonnegative remainder measures the deduction from valuing the deviation
at the initial marginal contribution, for expansion and contraction alike.

For a feasible replacement $z_i$, let $d_i=z_i-q_i$ and
$s_{i\ell}=k_\ell^\top d_i$.
Linearity of appeal gives $A_{i\ell}(z_i)=A_{i\ell}(q_i)+s_{i\ell}$.
Applying \eqref{eq:share-expansion} at every node yields
\begin{align}
  \mathcal M_i(z_i;r,B_i)-\mathcal M_i(q_i;r,B_i)
  &=\sum_\ell\frac{r_\ell B_{i\ell}}{D_\ell^2}k_\ell^\top d_i
  -\mathcal R_i^{\mathrm{sat}}(q,d_i),\label{eq:app-revenue-expansion}\\
  \mathcal R_i^{\mathrm{sat}}(q,d_i)
  &=\sum_\ell\frac{r_\ell B_{i\ell}(k_\ell^\top d_i)^2}
  {D_\ell^2(D_\ell+k_\ell^\top d_i)}\geq0.\label{eq:app-remainder}
\end{align}
Symmetry of $H_i$ gives the exact cost difference
\begin{equation}\label{eq:app-cost-expansion}
  C_i(q_i+d_i)-C_i(q_i)
  =(c_i+H_i q_i)^\top d_i+\tfrac12d_i^\top H_i d_i.
\end{equation}
Subtracting costs from margin gives, with the field in \eqref{eq:field},
\begin{equation}\label{eq:payoff-expansion}
  \pi_i(z_i,q_{-i})-\pi_i(q)
  =-F_i(q)^\top d_i-\tfrac12d_i^\top H_i d_i-\mathcal R_i^{\mathrm{sat}}(q,d_i).
\end{equation}
Thus a nonpositive first-order gain rules out a profitable deviation of any size.

For later use, consider a fraction $t\in[0,1]$ of a feasible deviation.
The new denominator is at least $b_\ell$, so
\begin{equation}\label{eq:app-remainder-order}
  0\leq\mathcal R_i^{\mathrm{sat}}(q,td_i)
  \leq t^2\sum_\ell\frac{r_\ell B_{i\ell}(k_\ell^\top d_i)^2}{D_\ell^2 b_\ell}.
\end{equation}
The right-hand coefficient is finite and independent of $t$.
This explicit bound justifies taking a directional limit at boundary capacities,
including profiles where some activities remain inactive.

\subsection{Equilibrium, uniqueness, and symmetry}\label{app:equilibrium}

The marginal incentive to expand is summarized by a vector field $F:\mathcal
Q\to\mathbb R^{Nm}$ whose $i$th block is the negative own-cash-flow gradient,
so that its coordinates are marginal expansion cost less marginal contribution:
\begin{equation}\label{eq:field}
  F_{ia}(q)=c_{ia}+(H_i q_i)_a-
  \sum_\ell r_\ell\frac{B_{i\ell}(q_{-i})}
  {[A_{i\ell}(q_i)+B_{i\ell}(q_{-i})]^2}K_{\ell a}.
\end{equation}
A negative coordinate is an incentive to add the corresponding capacity, and
at zero capacity a positive coordinate makes entry unprofitable.
At an equilibrium every activity with positive capacity has a zero
coordinate: its marginal contribution equals its marginal capacity cost.

Proposition~\ref{prop:equilibrium} follows by identifying Nash equilibria with
solutions of the variational inequality
\begin{equation}\label{eq:vi}
  \langle F(q^*),q-q^*\rangle\geq0\qquad\text{for every }q\in\mathcal Q,
\end{equation}
then establishing existence and strong monotonicity on the entire orthant.

\subsubsection{Own concavity and the Nash--VI equivalence}

Holding rivals fixed, the own margin Hessian is
\begin{equation}\label{eq:app-hessian}
  \nabla_i^2\mathcal M_i(q_i;r,B_i)
  =-2\sum_\ell\frac{r_\ell B_{i\ell}}{(A_{i\ell}+B_{i\ell})^3}
  k_\ell k_\ell^\top.
\end{equation}
Its quadratic form is nonpositive, and subtracting $H_i$ makes the
payoff Hessian
bounded above by $-\nu I$.
The formula is valid on an open neighborhood of each feasible profile
because all denominators
are strictly positive there.
Thus each own payoff is strictly concave on its convex feasible set,
with a unique
best response whenever a maximum exists.

Suppose $q^*$ solves the joint VI \eqref{eq:vi}.
Choose a feasible joint alternative that changes only firm $i$.
The VI then says $F_i(q^*)^\top(z_i-q_i^*)\geq0$.
The exact payoff expansion, positive cost curvature, and nonnegative
remainder imply
$\pi_i(z_i,q_{-i}^*)\leq\pi_i(q^*)$ for every feasible $z_i$.
This holds for every firm, so $q^*$ is Nash.
For the converse, suppose $q$ is Nash and fix $z_i\geq0$ with $d_i=z_i-q_i$.
The intermediate choice $q_i+td_i$ is feasible for $0<t\leq1$, and
Nash optimality gives
\begin{equation}\label{eq:app-nash-direction}
  0\geq-tF_i(q)^\top d_i-\tfrac12t^2d_i^\top H_i d_i-\mathcal
  R_i^{\mathrm{sat}}(q,td_i).
\end{equation}
Divide by $t$ and use \eqref{eq:app-remainder-order} as $t\downarrow0$.
It follows that $F_i(q)^\top d_i\geq0$.
Sum these inequalities across the blocks of an arbitrary feasible
joint alternative
to obtain \eqref{eq:vi}.
This proves necessity and sufficiency, including at the boundary of the orthant.

\subsubsection{Demand-side monotonicity}

Own concavity alone does not establish joint monotonicity: rival
responses also enter the joint field.
We therefore calculate its demand part explicitly.
At one node let $a=(a_1,\ldots,a_N)$ denote the nonnegative appeal vector and
$D=b+\sum_i a_i$.
The negative marginal-contribution field in appeal coordinates is
$f_i(a)=-r(D-a_i)/D^2$.
For a direction $u\in\mathbb R^N$ and $U=\sum_i u_i$, direct
differentiation gives
\begin{equation}\label{eq:app-node-derivative}
  (Df(a)u)_i=\frac r{D^3}\bigl[D u_i+(D-2a_i)U\bigr].
\end{equation}
The raw Jacobian need not be symmetric.
Its quadratic form, which equals that of its symmetric part, is
\begin{align}
  u^\top Df(a)u
  &=\frac r{D^3}\left[D\left(\sum_i u_i^2+U^2\right)-2U\sum_i a_i
  u_i\right]\label{eq:node-monotonicity}\\
  &=\frac r{D^3}\left[
    b\left(\sum_i u_i^2+U^2\right)
  +\sum_i a_i\left\{\sum_{j\ne i}u_j^2+(U-u_i)^2\right\}\right]\geq0.
  \label{eq:app-demand-squares}
\end{align}
The second equality expands $D=b+\sum_i a_i$ and completes a square
separately for each $a_i$.
All coefficients are nonnegative.
This is the property of the proportional-share technology that preserves
monotonicity in the interacting game.

For capacity profiles $q,v$, put $h=q-v$ and $q(t)=v+th$.
At node $\ell$, the corresponding appeal direction is
$u_{i\ell}=k_\ell^\top h_i$.
The margin part $F^R$ of the capacity field is the pullback of the
node fields through $K$.
The fundamental theorem of calculus therefore yields
\begin{equation}\label{eq:app-integrated-demand}
  \langle F^R(q)-F^R(v),h\rangle
  =\int_0^1\sum_\ell u_\ell^\top Df_\ell(A_\ell(q(t)))u_\ell\,\mathrm dt
  =\mathcal M(q,v)\geq0.
\end{equation}
The finite sums and continuous derivatives justify the integral.
A nonzero $h_i$ may lie in the null space of $K$, so this argument
need not supply
strict curvature in all capacity directions.
That strictness comes from the resource-cost matrices.
The node inequality is the unit-exponent case of the marginal-share Jacobian
curvature in \citet[Lemma~A.4]{ewerhart_quartieri_contests_2020}, with the
outside coordinate fixed and its directional change set to zero.
Equation~\eqref{eq:app-integrated-demand} transfers that curvature to capacity;
the cost bound controls directions in the null space of $K$, supplying the
strong monotonicity required by the general variational-inequality result
\citep{parise_ozdaglar_network_games_2019}.

\subsubsection{Gaussian cost curvature}

The Gaussian congestion kernel $K^c$ is positive semidefinite on any finite
set.
When $\tau_c=0$, it is the all-ones matrix and
$v^\top K^c v=(\sum_a v_a)^2\geq0$.
When $\tau_c>0$, factor its entries and expand the exponential:
\begin{equation}\label{eq:app-gaussian-series}
  K^c_{ab}=e^{-\tau_c\|x_a\|^2/d}e^{-\tau_c\|x_b\|^2/d}
  \sum_{n=0}^{\infty}\frac{(2\tau_c/d)^n}{n!}\langle x_a,x_b\rangle^n.
\end{equation}
For every nonnegative integer $n$, the matrix with entries
$\langle x_a,x_b\rangle^n$ is a Gram matrix of the tensor vectors
$x_a^{\otimes n}$;
for $n=0$ use the scalar vector one.
Consequently
\begin{equation}\label{eq:app-gaussian-psd}
  v^\top K^c v=\sum_{n=0}^{\infty}\frac{(2\tau_c/d)^n}{n!}
  \left\|\sum_a v_a e^{-\tau_c\|x_a\|^2/d}x_a^{\otimes n}\right\|^2\geq0.
\end{equation}
Finite activity sums and absolute convergence of the exponential
series justify the equality.
Thus $H=\eta K^c+\nu I$ has quadratic form at least $\nu\|v\|^2$ for $\eta\geq0$.
The same conclusion holds for any symmetric $H_i$ satisfying the stated bound.

\subsubsection{Existence on the unbounded cone}

The cost part of $F$ has difference $H_i(q_i-v_i)$ in block $i$.
Combining its curvature with \eqref{eq:app-integrated-demand} gives
\begin{equation}\label{eq:strong-monotonicity}
  \langle F(q)-F(v),q-v\rangle
  =\mathcal M(q,v)+\sum_i(q_i-v_i)^\top H_i(q_i-v_i)
  \geq\nu\|q-v\|^2.
\end{equation}
To establish existence, first obtain a strategy bound that holds
uniformly over rivals.
Each share lies between zero and one, hence
\begin{equation}\label{eq:app-capacity-bound}
  \pi_i(q_i,q_{-i})\leq R_\Sigma-c_i^\top q_i-\tfrac12q_i^\top H_i q_i
  \leq R_\Sigma-\tfrac\nu2\|q_i\|^2.
\end{equation}
Zero capacity earns zero for every rival profile.
Choose $R>\sqrt{2R_\Sigma/\nu}$ and let
$X_i=\{q_i\geq0:\|q_i\|\leq R\}$.
Every strategy with norm at least $R$ has negative payoff and is
dominated by zero.
The sets $X_i$ are nonempty, compact, and convex.
Continuity and strict own concavity give a unique maximizer over each
$X_i$ for every rival profile.

On the compact product $X=\prod_i X_i$, continuity of payoffs and uniqueness
of each maximizer make the best-response map continuous by the maximum theorem.
Brouwer's fixed-point theorem therefore supplies a fixed point in $X$.
No choice outside $X_i$ can improve on it, since those choices earn
negative payoff and zero is available inside.
The fixed point is therefore Nash on the entire orthant.

\subsubsection{Uniqueness and the symmetry boundary}\label{app:symmetry-proof}

Let $q^*$ and $v^*$ be two equilibria.
Their equivalent VI conditions, tested against one another, give
\begin{equation}\label{eq:app-uniqueness}
  \langle F(q^*),v^*-q^*\rangle\geq0,\qquad
  \langle F(v^*),q^*-v^*\rangle\geq0.
\end{equation}
Adding and reversing signs bounds the monotonicity pairing above by zero.
Its lower bound is $\nu\|q^*-v^*\|^2$, so the two equilibria coincide.
If primitives are invariant under relabeling firms, a permutation of
an equilibrium remains an equilibrium:
each permuted firm faces the same own cost and the correspondingly
permuted sum of rivals' appeals.
Uniqueness requires the profile to be invariant under every
transposition of firms.
Hence all firm blocks are equal.
Capability heterogeneity removes this permutation invariance and can
generate different equilibrium allocations.

\subsection{The no-interaction counterfactual}\label{app:independent-proof}

In the independent-firm economy, each firm faces only outside appeal while
retaining its own opportunity and cost primitives.
Its payoff is
\begin{equation}\label{eq:app-independent-payoff}
  \pi_i^{\mathrm M}(q_i)=\sum_\ell r_\ell\frac{k_\ell^\top
  q_i}{k_\ell^\top q_i+b_\ell}
  -c_i^\top q_i-\tfrac12q_i^\top H_i q_i.
\end{equation}
The associated field, derived by differentiation of this objective, is
\begin{equation}\label{eq:app-independent-field}
  F_i^{\mathrm M}(q_i)=c_i+H_i q_i-
  \sum_\ell\frac{r_\ell b_\ell}{(k_\ell^\top q_i+b_\ell)^2}k_\ell.
\end{equation}
With no rival blocks in its derivative, this operator is a Cartesian
product of one-firm proportional-share problems.
Each firm retains access to the full $R_\Sigma$, so aggregate market margin
can differ between the independent and strategic economies.

Each block Jacobian is
\begin{equation}\label{eq:app-independent-curvature}
  DF_i^{\mathrm M}(q_i)=H_i+
  2\sum_\ell\frac{r_\ell b_\ell}{(k_\ell^\top q_i+b_\ell)^3}k_\ell k_\ell^\top
  \succeq\nu I.
\end{equation}
Integration along each block segment and summation over firms gives
$\langle F^{\mathrm M}(q)-F^{\mathrm M}(v),q-v\rangle\geq\nu\|q-v\|^2$.
The payoff bound \eqref{eq:app-capacity-bound} still applies.
A maximum therefore exists for each block and strict concavity makes it unique;
the vector of these maxima is the unique independent equilibrium.
The exact payoff expansion remains valid with $B_{i\ell}=b_\ell$.
The VI equivalence and the distance and regret arguments below follow
directly for this product operator.
Under identical capabilities every independent optimum is identical,
even though its level need not equal
the symmetric strategic equilibrium.

\subsection{Response to opportunity and cost shifts}\label{app:response}

Throughout, $q$ is feasible, so $D_\ell\geq B_{i\ell}\geq b_\ell>0$ for every
firm and node.
Write $F(q;r)$ for the field \eqref{eq:field} at opportunity values $r$.

\subsubsection{Opportunity continuity}

Under Assumption~\ref{ass:primitives}, equilibria at nonnegative opportunity
vectors $r,r'$ satisfy
\begin{equation}\label{eq:response-continuity}
  \|q^*(r')-q^*(r)\|\leq\frac{L_r}{\nu}\|r'-r\|,
  \qquad L_r^2=N\sum_\ell\frac{\|k_\ell\|^2}{b_\ell^2}.
\end{equation}
This bound requires neither a fixed active set nor differentiability of the
equilibrium map.

Only the opportunity weights change, so the cost terms cancel and
\begin{equation}\label{eq:app-field-shift}
  F_{ia}(q;r')-F_{ia}(q;r)=-\sum_\ell(r'_\ell-r_\ell)\mathcal G_{ia,\ell},
  \qquad
  0\leq\mathcal G_{ia,\ell}=\frac{B_{i\ell}}{D_\ell^2}K_{\ell
  a}\leq\frac{K_{\ell a}}{b_\ell}.
\end{equation}
The upper bound uses $B_{i\ell}\leq D_\ell$ and $b_\ell\leq D_\ell$.
By Cauchy--Schwarz, each coordinate of the shift has square at most
$\|r'-r\|^2\sum_\ell K_{\ell a}^2/b_\ell^2$.
Summing over firms and activities gives
$\|F(q;r')-F(q;r)\|\leq L_r\|r'-r\|$ at every feasible $q$.
Let $q=q^*(r)$ and $q'=q^*(r')$.
Adding the two variational inequalities \eqref{eq:vi}, each evaluated
at the other
equilibrium, and using \eqref{eq:strong-monotonicity} at opportunity
values $r$ gives
\begin{equation}\label{eq:app-continuity}
  \nu\|q'-q\|^2\leq\langle F(q';r)-F(q';r'),q'-q\rangle
  \leq L_r\|r'-r\|\,\|q'-q\|,
\end{equation}
which is \eqref{eq:response-continuity}.
This is the strongly monotone sensitivity bound of
\citet{parise_ozdaglar_network_games_2019} with the model's loading bound.

\subsubsection{The Jacobian and its curvature}\label{app:response-jacobian}

Formally, write $k_\ell^\top$ for row $\ell$ of $K$ and $D_\ell=A_{i\ell}+B_{i\ell}$.
Raising $r_\ell$ lowers coordinate $(i,a)$ of the field \eqref{eq:field}
by exactly $\mathcal G_{ia,\ell}=B_{i\ell}K_{\ell a}/D_\ell^2$, so an
opportunity change $\delta r$ moves the field by $-\mathcal G\delta r$ at
fixed capacities.
The field's Jacobian $J$ has own and rival blocks
\begin{equation}\label{eq:jacobian}
  J_{ii}=H_i+\sum_\ell\frac{2r_\ell B_{i\ell}}{D_\ell^3}k_\ell k_\ell^\top,
  \qquad
  J_{ij}=\sum_\ell\frac{r_\ell(B_{i\ell}-A_{i\ell})}{D_\ell^3}k_\ell k_\ell^\top
  \quad(j\ne i).
\end{equation}
Let $\mathcal U$ be a fixed set of active firm-activity coordinates and $P$
the orthogonal projection onto them.
The projected linear response to an opportunity direction $\delta r$ is defined by
\begin{equation}\label{eq:linear-response}
  PJ\,\delta q=P\mathcal G\,\delta r,\qquad P\delta q=\delta q.
\end{equation}
Its $i$th block is~\eqref{eq:firm-response} on the active coordinates.
The direct response $\delta q^{\mathrm d}$ solves this system with $J$
replaced by $J_{\mathrm d}=\operatorname{blockdiag}(J_{ii})$.
It retains the strategic origin's rival fields $B_i$, rather than replacing
them by outside appeal alone.
The difference between the responses in Proposition~\ref{prop:response}
is the direct adjustment transmitted through the rival blocks
$J_{\mathrm o}=J-J_{\mathrm d}$.

Along a line $q+tu$ through a feasible profile, node $\ell$ contributes
$r_\ell(D_\ell-A_{i\ell})D_\ell^{-2}k_\ell$ to firm $i$'s block, with
numerator and
denominator affine in $t$.
Write $U_\ell=\sum_{j}k_\ell^\top u_j$ for the total appeal change and
$u_{i\ell}=k_\ell^\top u_i$ for firm $i$'s own change.
The quotient rule gives the derivative
\begin{equation}\label{eq:app-node-quotient}
  r_\ell\Bigl[\frac{U_\ell-u_{i\ell}}{D_\ell^2}
  -\frac{2(D_\ell-A_{i\ell})U_\ell}{D_\ell^3}\Bigr]k_\ell
  =-\frac{r_\ell}{D_\ell^3}\bigl[2B_{i\ell}u_{i\ell}
  +(B_{i\ell}-A_{i\ell})(U_\ell-u_{i\ell})\bigr]k_\ell,
\end{equation}
so the negative marginal-profit field has directional derivative $Ju$ with the
blocks in \eqref{eq:jacobian}.
The pairing $\langle Ju,u\rangle$ is therefore the derivative of
$t\mapsto\langle F(q+tu),u\rangle$ at zero.
That derivative is $\sum_{i}u_i^\top H_i u_i$ plus the node form
\eqref{eq:node-monotonicity}, which is nonnegative, so
$\langle Ju,u\rangle\geq\nu\|u\|^2$.
The own-block part contributes
$\sum_{i}u_i^\top H_i u_i+\sum_{i,\ell}2r_\ell B_{i\ell}D_\ell^{-3}u_{i\ell}^2$,
so $J_{\mathrm d}$ has the same modulus.
Neither bound uses an equilibrium or an active set.

\subsubsection{Existence, uniqueness, and size of the response}

Let $T$ be either $J$ or $J_{\mathrm d}$, and suppose $u$ solves
$PTu=Pg$ with $Pu=u$.
Because $P$ is an orthogonal projection,
\begin{equation}\label{eq:app-restricted-bound}
  \nu\|u\|^2\leq\langle Tu,u\rangle=\langle Tu,Pu\rangle
  =\langle PTu,u\rangle=\langle Pg,u\rangle\leq\|Pg\|\,\|u\|.
\end{equation}
Hence $\|u\|\leq\|Pg\|/\nu$.
Differences of solutions solve the system with zero forcing, so the
solution is unique.
The linear map $u\mapsto PTPu+(u-Pu)$ has trivial kernel by the same bound.
It is therefore onto in finite dimension, and a preimage of $Pg$ solves the
restricted system.
Thus at every feasible profile and for any fixed active set, both the full
and direct projected systems have unique solutions bounded by
$\|P\mathcal G\delta r\|/\nu$.
At an equilibrium with strict complementarity the derivative identification
below gives the economic interpretation in Proposition~\ref{prop:response}.

\subsubsection{Direct response and rival feedback}

Subtracting the direct system from the full system gives
$PJ(\chi\delta r-\delta q^{\mathrm d})=-PJ_{\mathrm o}\delta q^{\mathrm d}$.
Both directions lie in the active coordinate space, where $J$ is invertible
by \eqref{eq:app-restricted-bound}.
Applying that inverse gives \eqref{eq:feedback}, which proves the
common-origin comparison in Proposition~\ref{prop:response}.

\subsubsection{The derivative of the equilibrium
map}\label{app:equilibrium-derivative}

Suppose strict complementarity holds at $q=q^*(r)$: each coordinate has positive
capacity or a strictly positive field value.
With $\mathcal U$ the coordinates of positive capacity,
\begin{equation}\label{eq:equilibrium-derivative}
  \|q^*(r')-q^*(r)-\chi\,(r'-r)\|=o(\|r'-r\|),
\end{equation}
so $\chi$ is the derivative of the equilibrium map and inactive
coordinates do not
respond.

The field is a rational function of capacities with positive
denominators, so it is
Fréchet differentiable, and its derivative along every line is $J$
(Appendix~\ref{app:response-jacobian}).
At an equilibrium every field coordinate is nonnegative and vanishes
where capacity
is positive.
Let $m>0$ bound from below each positive capacity and
each positive field value at $q$.
By \eqref{eq:response-continuity} and the Lipschitz continuity of $F$
on the orthant,
for $\|r'-r\|$ small the equilibrium $q'=q^*(r')$ keeps every
positive coordinate
positive, and every coordinate with a positive field value stays at zero.
Both equilibria therefore satisfy $PF=0$ with $P(q'-q)=q'-q$.
Expanding, $F(q';r')-F(q;r)=J(q'-q)-\mathcal G\,\delta r+R$, where
$R$ collects the
Fréchet remainder and the change in $\mathcal G$ between $q$ and $q'$; both are
$o(\|q'-q\|+\|\delta r\|)$.
The error $q'-q-\chi\delta r$ is supported on $\mathcal U$ and solves
the linearized
system with forcing $-R$, so \eqref{eq:app-restricted-bound} bounds it by
$\|PR\|/\nu=o(\|\delta r\|)$.
This proves \eqref{eq:equilibrium-derivative}.

Linear costs enter the field additively, $F(q;c')-F(q;c)=c'-c$.
The same argument with forcing $-(c'-c)$ in place of $\mathcal
G\,\delta r$ shows that,
under strict complementarity and for nearby nonnegative costs $c'$,
\begin{equation}\label{eq:cost-response}
  \|q^*(c')-q^*(c)-u\|=o(\|c'-c\|),\qquad PJu=-P(c'-c),\quad Pu=u.
\end{equation}
The forcing does not depend on capacities, so the remainder is the
Fréchet remainder
alone.
If $c'-c$ is supported on one firm, every other firm's problem
against frozen rivals
is unchanged, so its frozen-rival response is zero.
The opportunity and cost responses are established separately, each for one
type of disturbance at a time.

\subsection{Proof of the composition-feedback
benchmark}\label{app:composition-feedback}

For $S=\mathbf1^\top q>0$, the quotient rule gives
$Dp_q(h)=(h-p\mathbf1^\top h)/S$.
Its kernel consists exactly of $h=\lambda q$: substituting this form gives
zero, and a zero numerator gives $h=(\mathbf1^\top h/S)q$.
Linearity yields \eqref{eq:proportional-feedback}.

At the affected activity the negative marginal payoff is $g_i=k_i+\omega_i
A_i-rB_i/D^2$, where $k_i$ is its linear cost and $\omega_i>0$ its curvature.
Differentiating the interior first-order conditions gives
\begin{equation}\label{eq:separable-jacobian}
  J=\operatorname{diag}(a)+c\mathbf1^\top,\qquad
  a_i=\omega_i+r/D^2>0,\qquad c_i=r(B_i-A_i)/D^3>0.
\end{equation}
Indeed the own derivative is $\omega_i+2rB_i/D^3=a_i+c_i$ and every rival
derivative in row $i$ is $c_i$.
Let $u$ denote the capacity derivative and $T=\sum_i u_i$.
The cost reduction imposes $a_i u_i+c_i T=s\,\mathbf{1}_{\{i=j\}}$.
Dividing by $a_i$ and summing yields
\begin{equation}\label{eq:separable-response}
  \begin{aligned}
    \Delta&=1+\sum_i c_i/a_i>0,&
    T&=\frac{s/a_j}{\Delta}>0,\\
    u_j&=\frac{s}{a_j}\frac{1+\sum_{i\ne j}c_i/a_i}{\Delta}>0,&
    u_i&=-\frac{c_i}{a_i}T<0\quad(i\ne j).
  \end{aligned}
\end{equation}
For uniqueness, the difference of two solutions has aggregate $U$ satisfying
$\Delta U=0$, hence each difference coordinate is zero.
Each unaffected activity has zero forcing and an invertible block by
positive-curvature strong monotonicity, so its response vanishes.
If firm $i$ has positive capacity $q_{ib}$ at another activity $b$, its share
derivative there is $-q_{ib}u_i/S_i^2\ne0$.
This proves the business-mix claim under the maintained interior path.
Entry and exit require a directional analysis; these derivatives make no
assertion across active-set changes.

\subsection{Two-site exposure ordering}\label{app:two-site-example}

Whether a capacity response separates firms' business mixes depends on
their initial exposures.
To make that dependence explicit, consider two firms at two distinct activity
sites with identity demand loading and positive diagonal quadratic cost
curvature.
Initial interior capacities are $q_i=(q_{i1},q_{i2})$ and
$q_j=(q_{j1},q_{j2})$, all positive.
At activity 1, the node-1 opportunity value and outside appeal satisfy $r,b>0$,
and each firm's payoff component has the form
\begin{equation}\label{eq:sign-payoff}
  \pi_i^{(1)}(q_{i1},q_{j1})
  =\frac{rq_{i1}}{b+q_{i1}+q_{j1}}-c_i q_{i1}-\frac{\omega_i}{2}q_{i1}^2,
  \qquad \omega_i,\omega_j>0.
\end{equation}
Firm $j$ receives a small cost reduction $c_j\mapsto c_j-s$, preserving
nonnegative costs.
Assume a differentiable interior equilibrium path satisfying both first-order
conditions locally, with activity 2 unchanged.
To compare business mixes we use the quadratic Wasserstein distance $W_2$,
the least root-mean-square distance over which mass must move to turn one
composition into the other.
The two sites are a distance $d_*>0$ apart.
Define $P_k(s)$ from the normalized two-site capacities and
$Y^F_{ij}(s)=W_2^2(P_i(s),P_j(0))-W_2^2(P_i(0),P_j(0))$, the change in the
squared distance from firm $i$'s business mix to firm $j$'s initial mix.

\begin{proposition}[Conditional capacity and composition
  response]\label{prop:two-site-adjustment}
  In the two-site economy and along the path just specified, suppose
  \begin{equation}\label{eq:sign-conditions}
    q_{i1}<b+q_{j1},\qquad q_{j1}q_{i2}>q_{i1}q_{j2}.
  \end{equation}
  Then the cost recipient expands in activity 1, $\dot q_{j1}>0$, its competitor
  contracts, $\dot q_{i1}<0$, and $\dot Y^F_{ij}>0$.
  Reversing the strict initial exposure ordering reverses the sign of
  $\dot Y^F_{ij}$ while preserving the capacity-response signs.
\end{proposition}

\subsubsection{Proof of the two-site capacity and composition
response}\label{app:two-site-adjustment}

\begin{proof}[Proof of Proposition~\ref{prop:two-site-adjustment}]
  Differentiate the two interior first-order conditions in
  \eqref{eq:sign-payoff}.
  Write $D=b+q_{i1}+q_{j1}$, $B_i=b+q_{j1}$, $B_j=b+q_{i1}$, and define
  \begin{equation}\label{eq:sign-jacobian}
    \begin{aligned}
      J_{ii}&=\omega_i+2rB_i/D^3,& J_{jj}&=\omega_j+2rB_j/D^3,\\
      J_{ij}&=r(B_i-q_{i1})/D^3,& J_{ji}&=r(B_j-q_{j1})/D^3,\\
      \Delta&=J_{ii}J_{jj}-J_{ij}J_{ji}.
    \end{aligned}
  \end{equation}
  These are the own and rival entries of the field's
  Jacobian~\eqref{eq:jacobian} restricted to activity 1.
  The demand-only part of the determinant has numerator
  $r^2[4{B_i}B_j-(B_i-q_{i1})(B_j-q_{j1})]$ over $D^6$.
  The bracket equals $3b^2+4b(q_{i1}+q_{j1})+(q_{i1}+q_{j1})^2>0$.
  The additional terms from $\omega_i,\omega_j>0$ are positive, so $\Delta>0$.
  Solving the differentiated conditions gives
  \begin{equation}\label{eq:sign-capacity}
    \begin{pmatrix}J_{ii}&J_{ij}\\J_{ji}&J_{jj}
    \end{pmatrix}
    \begin{pmatrix}\dot q_{i1}\\\dot q_{j1}
    \end{pmatrix}
    =
    \begin{pmatrix}0\\1
    \end{pmatrix},
    \qquad \dot q_{i1}=-J_{ij}/\Delta,\quad \dot q_{j1}=J_{ii}/\Delta.
  \end{equation}
  The substitution inequality $q_{i1}<b+q_{j1}$ makes $J_{ij}>0$, yielding
  competitor contraction and recipient expansion.

  For shares $p,p'\in[0,1]$ on two sites separated by $d_*>0$, the exact
  squared transport distance is
  \begin{equation}\label{eq:two-site-transport}
    W_2^2\bigl((p,1-p),(p',1-p')\bigr)=d_*^2|p-p'|.
  \end{equation}
  Any feasible nonnegative coupling has off-diagonal masses whose
  difference is $p-p'$.
  Their sum is at least $|p-p'|$, and matching the common mass at each site
  attains the bound.
  Put $p_k=q_{k1}/(q_{k1}+q_{k2})$.
  Because activity 2 is unchanged,
  $\dot p_i=q_{i2}\dot q_{i1}/(q_{i1}+q_{i2})^2<0$.
  The exposure ordering implies $p_j(0)>p_i(0)$ and persists locally against
  the frozen $p_j(0)$.
  Differentiating the squared distance therefore gives
  \begin{equation}\label{eq:sign-distance}
    \dot Y^F_{ij}=-d_*^2\dot p_i
    =\frac{d_*^2{q_{i2}}J_{ij}}{(q_{i1}+q_{i2})^2\Delta}>0.
  \end{equation}
  If $p_j(0)<p_i(0)$, the derivative instead equals $d_*^2\dot p_i<0$.
  This establishes the statement and its reversed-order conclusion.
\end{proof}

Initial exposure ordering determines whether contraction moves the
competitor's business mix toward or away from the recipient's initial mix;
the capacity-response signs alone do not determine that direction.


\subsection{Cash-flow exposure and competitive news}\label{app:cash-flow-support}

\subsubsection{Cash-flow exposure}

The node shares also determine operating cash-flow exposure.
Write $X(q)$ for the matrix of node shares
$X_{i\ell}=A_{i\ell}/D_\ell$.

\begin{corollary}[Cash-flow exposure]\label{cor:exposure}
  At fixed capacities, payoff is affine in opportunity values,
  $\pi_i(q;r')-\pi_i(q;r)=\sum_\ell(r'_\ell-r_\ell)X_{i\ell}(q)$.
  If opportunity innovations $\varepsilon$ have covariance $\Omega$, then
  \begin{equation}\label{eq:exposure-covariance}
    \operatorname{Cov}\bigl(X(q)\varepsilon\bigr)=X(q)\,\Omega\,X(q)^\top.
  \end{equation}
  A rival's expansion weakly lowers the firm's payoff: if $q_j'\geq q_j$ for
  some $j\neq i$, then
  \begin{equation}\label{eq:rival-cross}
    \pi_i(q_j',q_{-j};r)\leq\pi_i(q;r).
  \end{equation}
\end{corollary}

Uncorrelated innovations therefore already produce correlated operating
cash flow whenever firms' appeal overlaps at common nodes.
Any strategic contribution to cash-flow covariance has to be measured
against this shared-exposure baseline.

\begin{proof}
  At fixed capacities, the margin is linear in $r$ with coefficients
  $A_{i\ell}/D_\ell$, and costs do not depend on $r$.
  The covariance formula \eqref{eq:exposure-covariance} is bilinearity of
  covariance for square-integrable innovations and a predetermined $X(q)$.
  If $q_j'\geq q_j$ for a rival $j$, then $A_{i\ell}$ and $C_i$ are unchanged
  while each $D_\ell$ weakly rises because $K\geq0$.
  Each term $r_\ell A_{i\ell}/D_\ell$ weakly falls, which gives
  \eqref{eq:rival-cross}.
\end{proof}

\begin{remark}[Envelope split]\label{rem:envelope}
  Along a direction $(u,\delta r)$ of capacities and opportunity
  values, firm $i$'s
  cash flow changes at the rate
  \begin{equation}\label{eq:envelope}
    \sum_\ell\delta r_\ell X_{i\ell}-\sum_a F_{ia}u_{ia}
    -\sum_\ell\frac{r_\ell A_{i\ell}}{D_\ell^2}\Bigl(\sum_{j}k_\ell^\top
    u_j-k_\ell^\top u_i\Bigr).
  \end{equation}
  At a static equilibrium $F_{ia}=0$ wherever the response moves the firm's
  capacity, so the change is direct opportunity news plus the value effect of
  rivals' movement.
  Along the equilibrium response $u=\chi\delta r$, under strict complementarity,
  the equilibrium cash-flow sensitivity is therefore
  given by \eqref{eq:equilibrium-cash-flow}:
  direct opportunity news plus the value of rivals' equilibrium repositioning.
\end{remark}

\begin{proof}
  The margin is linear in $r$, so the opportunity term is the share $X_{i\ell}$.
  The quotient rule gives each node's capacity derivative, and symmetry of $H_i$
  collects the own terms into $-F_i^\top u_i$.
  The remaining node terms involve only rivals' appeal changes.
  For \eqref{eq:equilibrium-cash-flow}, cash flow is Fréchet differentiable and
  shares are continuous, so \eqref{eq:equilibrium-derivative} replaces
  $q^*(r+t\delta r)-q^*(r)$ by $t\chi\delta r$ up to $o(t)$.
  The own term vanishes because $\chi\delta r$ is supported where $F$ vanishes.
\end{proof}

\subsubsection{Common and differential competitive news}\label{app:common-news}

The retained share in Proposition~\ref{prop:common-news} measures the
common-node component, separately from its own-appeal correction.
Let $\vartheta_i=\Lambda_i\psi_i$, with mean
$\bar\vartheta=N^{-1}\sum_i\vartheta_i$, and let $\zeta_i=-\vartheta_i^\top g$ be
the common-node part of~\eqref{eq:rival-return}, with cross-sectional mean
$\bar\zeta$.
Across disturbances, treat node growth $g$ as a random vector with
second-moment matrix $\Sigma_g=\mathbb E[gg^\top]$.
Whenever $\sum_i\vartheta_i^\top\Sigma_g\vartheta_i>0$, the fraction of the
common-node part's expected cross-sectional second moment that survives
removal of the cross-sectional mean is
\begin{equation}\label{eq:retained-share}
  \varrho=\frac{\mathbb E\sum_i(\zeta_i-\bar\zeta)^2}{\mathbb E\sum_i\zeta_i^2}
  =\frac{\sum_i(\vartheta_i-\bar\vartheta)^\top\Sigma_g(\vartheta_i-\bar\vartheta)}
  {\sum_i\vartheta_i^\top\Sigma_g\vartheta_i}.
\end{equation}
The common-node projection identity is algebraic; it does not establish
that demeaning by itself makes competitive news statistically undetectable.

\begin{proof}[Proof of Proposition~\ref{prop:common-news}]
  By \eqref{eq:envelope}, the cash-flow change along $(u,\delta r)$ is direct
  opportunity news, an own term $-\sum_a F_{ia}u_{ia}$, and the rival term
  $-\sum_\ell r_\ell A_{i\ell}D_\ell^{-2}(\sum_j k_\ell^\top
  u_j-k_\ell^\top u_i)$.
  The own term vanishes because $u$ is supported where $F$ vanishes.
  A linear-cost direction $\delta c$ adds the direct effect
  $-\delta c_i^\top q_i$, which does not involve rivals, so the rival
  term is unchanged.
  Since $r_\ell A_{i\ell}/D_\ell=r_\ell X_{i\ell}=\mathcal M_i\psi_{i\ell}$ and
  $\sum_j k_\ell^\top u_j-k_\ell^\top u_i=D_\ell g_\ell-k_\ell^\top u_i$,
  dividing the rival term by $\pi_i$ gives \eqref{eq:rival-return}.
  For the retained fraction,
  $\sum_i(\vartheta_i^\top g-\bar\vartheta^\top
  g)^2=\sum_i((\vartheta_i-\bar\vartheta)^\top g)^2$,
  whose expectation is
  $\sum_i(\vartheta_i-\bar\vartheta)^\top\Sigma_g(\vartheta_i-\bar\vartheta)$;
  the expected second moment before demeaning is
  $\sum_i\vartheta_i^\top\Sigma_g\vartheta_i$.
  If all $\vartheta_i$ are equal, every deviation is zero.
\end{proof}

\subsection{Numerical equilibrium and composition
certificates}\label{app:certificates}

Projected extragradient locates each equilibrium; independently evaluated
residuals determine whether to accept it.
For a feasible numerical profile $\widehat q$, define the orthant
stationarity residual by
\begin{equation}\label{eq:residual}
  s_{ia}(\widehat q)=
  \begin{cases}
    F_{ia}(\widehat q),&\widehat q_{ia}>0,\\
    \min\{F_{ia}(\widehat q),0\},&\widehat q_{ia}=0.
  \end{cases}
\end{equation}
At an inactive coordinate, positive $F_{ia}$ discourages entry and contributes
no error; a negative value records an unexploited incentive to enter.

\subsubsection{Capacity error and unilateral regret}\label{app:regret-proof}

\begin{proposition}[Equilibrium accuracy and unilateral payoff
  regret]\label{prop:certificate}
  Under Assumption~\ref{ass:primitives}, the strategic and
  independent operators each satisfy
  \begin{equation}\label{eq:certificate}
    \|\widehat q-q^*\|\leq\frac{\|s(\widehat q)\|}{\nu},\qquad
    \operatorname{Regret}_i(\widehat q)\leq\frac{\|s_i(\widehat q)\|^2}{2\nu},
  \end{equation}
  where regret is the supremum of the unilateral payoff improvement
  holding rivals fixed.
\end{proposition}

\begin{proof}
  Put $n=F(\widehat q)-s$.
  By construction $n\geq0$ and $n_{ia}\widehat q_{ia}=0$.
  For every feasible $z$,
  \begin{equation}\label{eq:app-approx-vi}
    \langle F(\widehat q),z-\widehat q\rangle
    =\langle s,z-\widehat q\rangle+\langle n,z\rangle
    \geq-\|s\|\|z-\widehat q\|.
  \end{equation}
  With $h=\widehat q-q^*$, this inequality at $z=q^*$ and the equilibrium VI
  give
  \begin{equation}\label{eq:app-distance-chain}
    \nu\|h\|^2\leq\langle F(\widehat q)-F(q^*),h\rangle
    \leq\|s\|\|h\|.
  \end{equation}
  Dividing by $\nu\|h\|$ when $h\ne0$ proves the distance bound; the case
  $h=0$ is immediate.
  \label{app:independent-distance}
  The independent product operator has the same modulus, so the bound applies
  to its unique equilibrium as well.

  For a feasible unilateral deviation $d_i=z_i-\widehat q_i$, the block form
  of \eqref{eq:app-approx-vi} and the exact payoff expansion give
  \begin{align}
    \pi_i(z_i,\widehat q_{-i})-\pi_i(\widehat q)
    &\leq\|s_i\|\|d_i\|-\tfrac\nu2\|d_i\|^2\nonumber\\
    &=\frac{\|s_i\|^2}{2\nu}
    -\frac\nu2\left(\|d_i\|-\frac{\|s_i\|}{\nu}\right)^2
    \leq\frac{\|s_i\|^2}{2\nu}.\label{eq:app-regret-square}
  \end{align}
  Taking the supremum over $z_i\geq0$ proves the regret bound.
  \label{app:independent-regret}
  The independent payoff expansion has the same form, giving the identical
  bound for each independent firm's objective.
\end{proof}

\subsubsection{Accuracy after capacity normalization}\label{app:observation}

Because composition divides capacities by total mass, a small capacity error
can matter when that mass is small.
For feasible profiles, the product norm bounds each block,
\begin{equation}\label{eq:app-block-bound}
  \|\widehat q_i-q_i^*\|^2\leq\sum_j\|\widehat q_j-q_j^*\|^2
  =\|\widehat q-q^*\|^2,
\end{equation}
so the joint certificate supplies a conservative error radius for every firm.

\begin{proposition}[Observation error under
  normalization]\label{prop:normalization}
  Suppose $\|\widehat q_i-q_i^*\|\leq e$ and
  $\widehat S_i=\mathbf1^\top\widehat q_i>\sqrt m\,e$.
  Then $S_i^*>0$ and
  \begin{equation}\label{eq:normalization-bound}
    \left\|\frac{\widehat q_i}{\widehat S_i}-\frac{q_i^*}{S_i^*}\right\|
    \leq\frac{2\sqrt m\,e}{\widehat S_i-\sqrt m\,e}.
  \end{equation}
\end{proposition}

\begin{proof}
  Cauchy--Schwarz bounds the mass error by
  \begin{equation}\label{eq:app-mass-error}
    |\widehat S_i-S_i^*|
    =|\langle\mathbf1,\widehat q_i-q_i^*\rangle|\leq\sqrt m\,e.
  \end{equation}
  Hence $S_i^*\geq\widehat S_i-\sqrt m\,e>0$.
  Nonnegativity gives $\|\widehat q_i\|\leq\widehat S_i$, and the identity
  \begin{equation}\label{eq:app-normalization-split}
    \frac{\widehat q_i}{\widehat S_i}-\frac{q_i^*}{S_i^*}
    =\frac{\widehat q_i-q_i^*}{S_i^*}
    +\frac{S_i^*-\widehat S_i}{\widehat S_i S_i^*}\widehat q_i
  \end{equation}
  therefore implies
  \begin{equation}\label{eq:app-normalization-chain}
    \left\|\frac{\widehat q_i}{\widehat S_i}-\frac{q_i^*}{S_i^*}\right\|
    \leq\frac{(1+\sqrt m)e}{S_i^*}
    \leq\frac{2\sqrt m\,e}{\widehat S_i-\sqrt m\,e},
  \end{equation}
  where the last step uses $m\geq1$.
\end{proof}

The benchmark experiments accept profiles using the independently evaluated
relative stationarity-distance target \ppnume[0]{ir-solve-target}.
If multiple starts give error radii $e_a,e_b$, uniqueness and the triangle
inequality bound their distance by $e_a+e_b$; agreement across starts
supplements the residual check.
These inequalities hold in real arithmetic, while their numerical evaluations
use floating-point capacities and residuals without interval-arithmetic
enclosures.
